\documentclass[10pt,a4paper]{amsart}
\usepackage[margin=19mm]{geometry}
\usepackage{amsmath,amssymb,mathtools}
\usepackage[scr=rsfs,cal=euler]{mathalfa}
\usepackage{xfrac}
\usepackage[unicode,hidelinks,bookmarks=false]{hyperref}
\newcommand{\ie}{i.e.\ }
\newcommand{\cf}{cf.\ }
\newcommand{\resp}{respectively\ }
\newcommand{\Subsection}{\subsection}
\providecommand{\nobreakdash}{\nobreak\hbox{-}}
\setlength{\emergencystretch}{2em}
\multlinegap0pt
\usepackage{amssymb}
\usepackage{enumerate}
\usepackage{tikz}
\newtheorem{theo}{Theorem}
\newtheorem{lem}{Lemma}
\title{\LARGE M\Large \MakeLowercase{eromorphic Functions and Linearization Phenomena in Partial Differential Equations}}
\author{Sujoy Majumder, Debabrata Pramanik, Jhilik Banerjee}
\date{Source: arXiv:2605.07636v1 (CC BY 4.0)}
\begin{document}
\maketitle

\noindent\emph{Source and licence.} \LARGE M\Large \MakeLowercase{eromorphic Functions and Linearization Phenomena in Partial Differential Equations}. Version arXiv:2605.07636v1 (CC BY 4.0), distributed by arXiv under the Creative Commons Attribution 4.0 International licence, http://creativecommons.org/licenses/by/4.0/, as recorded in the arXiv licence metadata for this exact version and shown on the abstract page https://arxiv.org/abs/2605.07636v1. Full licence notice: \texttt{licenses/arxiv-2605.07636-CC-BY-4.0.txt}.\par\smallskip

\noindent\textit{Editorial scope.} Counted unit: the article's lem L2 (source0.tex lines 441--446). The article's complete original proof of it is delivered with it (source0.tex lines 461--661, one \texttt{proof} environment of 201 lines, which the article places away from the statement). No mathematics is written, repaired or re-derived: the statement and the proof are the article's own lines, in the article's own notation, and no retained line is substituted or re-flowed.  Retained uncounted as the source-local context the proof needs: lines 141--143; lines 154--156; lines 162--172; lines 344--348; lines 395--399; lines 448--455.  Nothing was omitted from any retained span, so the package carries no omission record.  The retained lines cite 1 works, whose entries are transcribed verbatim from the article's own bibliography source in the e-print and delivered with short editorial labels recorded in the item's reference mapping.  No retained line contains a figure environment, an image-inclusion command or drawing code, so no image is delivered and the package declares no asset dependency.  Editorial formatting only, in addition: an AMS \texttt{amsart} preamble that restates the article's own declarations and macros used by the retained lines, namely the environments \texttt{theo}, \texttt{lem} on separate counters, together with the standard packages those lines need.  The retained environments print in this document as Theo 1, Lemma 1 in order of appearance, whereas the article prints the same items with the numbering of its own full document.  The complete native source of the article as distributed in the arXiv e-print for this exact version is bundled unchanged as \texttt{sources/200100/source0.tex}.
\begin{align}\label{eq}
G^g_f(z)=f(g(z),g(z),\ldots,g(z)).
\end{align}

\begin{align}\label{Eq}
\frac{\partial h(z)}{\partial z_i}=a G^g_{h}(z)+bh(z)+c,
\end{align}

\begin{theo}\label{t1.1} Suppose that $f$ is a non-constant meromorphic function in $\mathbb{C}$ and $g$ is an entire function in $\mathbb{C}^n$ satisfying the equation (\ref{Eq}),
where $z\in\mathbb{C}^n$, $i\in\{1,2,\ldots,n\}$, $a\neq 0, b,c$ are constants in $\mathbb{C}$, $h(z)=f(z_1+z_2+\ldots+z_n)$ and $G^g_h(z)$ is defined by (\ref{eq}). Then 
\begin{enumerate}
\item[(i)] $g$ must be linear in $\mathbb{C}^n$, if $f$ is transcendental;
\item[(ii)] $g$ must be a polynomial in $\mathbb{C}^n$ of degree less than or equal to $2$, if $f$ is a rational function in $\mathbb{C}$. Furthermore the degree of $g$ is $2$ if and only if 
\begin{align*}
f(z)=\beta+\frac{\alpha}{z-w_0},\quad ng(z)=w_0-a\left(z_1+\ldots+z_n-w_0\right)^2
\end{align*}
and $b=a\beta+c=0$, where $\alpha\neq 0, \beta, w_0$ are complex numbers.
\end{enumerate}
\end{theo}

\begin{lem}\label{Lm2} Let $f$ be a meromorphic function in $\mathbb{C}$ and let $g:\mathbb{C}^n\to \mathbb{C}$ be an entire function. Suppose $f$ and $g$ are transcendental. Then 
\begin{align*}
\limsup\limits_{r\to\infty}\frac{T(r,f(g))}{T(r,f)}=\infty.
\end{align*}
\end{lem}

\begin{lem}\label{L1.1} Let $f:\mathbb{C}\to \mathbb{P}^1$ be a meromorphic function. If $g:\mathbb{C}^n\to \mathbb{C}$ is a non-constant polynomial of degree $m$, then 
	\begin{align*}
		T(r,f(g))\geq (1-o(1))T(r^m,f).
	\end{align*}
	\end{lem}

\begin{lem}\label{L3} \cite[Theorem 1.26]{Hu-Li-Yang-2003} Let $f:\mathbb{C}^n\to\mathbb{P}^1$ be a non-constant meromorphic function. Assume that 
$R(z, w)=\frac{A(z, w)}{B(z, w)}$. Then
\begin{align*}
 T\left(r, R_f\right)=\max \{p, q\} T(r, f)+O\Big(\sideset{}{_{j=0}^{p}}{\sum}T(r, a_j)+\sideset{}{_{j=0}^{q}}{\sum}T(r, b_j)\Big),
 \end{align*}
for all large $r$ outside a set $E$ with $\int_E d\log r<\infty$, where $R_f(z)=R(z, f(z))$ and two coprime polynomials $A(z, w)$ and $B(z,w)$ are given
respectively $A(z,w)=\sum_{j=0}^p a_j(z)w^j$ and $B(z,w)=\sum_{j=0}^q b_j(z)w^j$.
\end{lem}

\begin{lem}\label{L2} \cite[Lemma 1.37]{Hu-Li-Yang-2003} Let $f:\mathbb{C}^n\to\mathbb{P}^1$ be a non-constant meromorphic function and $I=(\alpha_1,\ldots,\alpha_n)\in \mathbb{Z}^n_+$ be a multi-index. Then for any $\varepsilon>0$, we have
\begin{align*}
m\left(r,\frac{\partial^I(f)}{f}\right)\leq |I|\log^+T(r,f)+|I|(1+\varepsilon)\log^+\log T(r,f)+O(1)
\end{align*}
for all large $r$ outside a set $E$ with $\int_E d\log r<\infty$.
\end{lem}

\begin{proof} From (\ref{Eq}), we have
\begin{align}\label{EqT1.0}
\frac{\partial h(z)}{\partial z_i}=aG^g_h(z)+bh(z)+c
\end{align}
where $z=(z_1,z_2,\ldots,z_n)$ and $i\in\{1,2,\ldots,n\}$. 
Note that
\begin{align}\label{EqT1.1}
N\left(r,\frac{\partial h(z)}{\partial z_i}\right)\leq 2N(r,h(z)),
\end{align}
where $z=(z_1,z_2,\ldots,z_n)$ and $i\in\{1,2,\ldots,n\}$.
Now using Lemma \ref{L2} and (\ref{EqT1.1}) to (\ref{EqT1.0}), we deduce that
\begin{align}\label{EqT1.2}
T\left(r,\frac{\partial h(z)}{\partial z_i}-bh(z)\right)&=N\left(r,\frac{\partial h(z)}{\partial z_i}-bh(z)\right)+m\left(r,\frac{\partial h(z)}{\partial z_i}-hf(z)\right)\\&\leq 
2N(r,h(z))+m\left(r,\frac{\frac{\partial h(z)}{\partial z_i}-bh(z)}{h(z)}\right)+m(r,h)\nonumber\\&\leq
2T(r,h)+o(T(r,h))\nonumber
\end{align}
outside possibly a set of finite logarithmic measure, where $i\in\{1,2,\ldots,n\}$.

Therefore using Lemma \ref{L3} and (\ref{EqT1.2}) to (\ref{EqT1.0}), we obtain
\begin{align}\label{EqT1.3}
T(r,G^g_h)=T\left(r, \frac{\frac{\partial h(z)}{\partial z_i}-bh(z)}{a}-\frac{c}{a}\right)&\leq T\left(r,\frac{\partial h(z)}{\partial z_i}-bh(z)\right)+O(1)\\&\leq 
2T(r,h)+o(T(r,h))\nonumber
\end{align}
outside possibly a set of finite logarithmic measure, where $i\in\{1,2,\ldots,n\}$.


Now we consider the following two cases.\par

\medskip
{\bf Case 1.} Let $g$ be a transcendental entire function in $\mathbb{C}^n$. Now from (\ref{EqT1.0}), we have 
\begin{align}\label{EqT1.4}
\frac{\partial h(z)}{\partial z_i}-bh(z)-c=aG^g_h(z).
\end{align}


\medskip
First we suppose that $f$ is a rational function in $\mathbb{C}$. Clearly $h$ is a rational function in $\mathbb{C}^n$. Let $h=\frac{P}{Q}$, where $P$ and $Q$ are rational functions in $\mathbb{C}^n$ of degree $p$ and $q$ respectively. Then by Lemma \ref{L3}, we have 
\begin{align*}
T(r, G^g_h)=\max\{p,q\}\;T(r,g)+O(1)\neq O(\log r),
\end{align*}
which shows that $G^g_h$ is transcendental. It is easy to conclude from (\ref{EqT1.4}) that $h$ is a rational function in $\mathbb{C}^n$ and so $f$ is a rational function in $\mathbb{C}$, which is impossible.


\medskip
Next we suppose that $f$ is a transcendental meromorphic function in $\mathbb{C}$. Consequently $h$ is also a transcendental meromorphic function in $\mathbb{C}^n$. Then from (\ref{EqT1.3}), we obtain 
\begin{align*}
\limsup\limits_{r\not\in E, r\to\infty}\frac{T\left(r,G^g_h\right)}{T(r,h)}<\infty,
\end{align*}
which is impossible by Lemma \ref{Lm2}.


\medskip
{\bf Case 2.} Let $g$ be a polynomial in $\mathbb{C}^n$. Suppose $\deg(g)=m$. If $m =1$, then
the conclusions $(i)$ and $(ii)$ of Theorem \ref{t1.1} both hold. Next suppose $m\geq 2$. Then by Lemma \ref{L1.1}, we get
\begin{align}\label{EqT1.5}
T\left(r,G^g_h\right)\geq (1-\varepsilon_1)T(r^m,h)
\end{align}
for any $\varepsilon_1>0$. We know that $T(r, h)$ is a convex function of $\log r$. Thus, for large $r$, we have 
\begin{align*}
\frac{T(r,h)-T(1,h)}{\log r-\log 1}\leq \frac{T(r^m,h)-T(1,h)}{\log r^m-\log 1},
\end{align*}
which shows that 
\begin{align}\label{EqT1.6}
T(r,h)\leq \frac{1}{m}(1+\varepsilon_2)T\left(r^m,h\right)
\end{align}
for any $\varepsilon_2>0$ and large $r$. Now using (\ref{EqT1.3}), (\ref{EqT1.5}) to (\ref{EqT1.6}), we obtain
\begin{align}\label{EqT1.7}
mT(r,h)\leq \frac{1+\varepsilon_2}{1-\varepsilon_1}\left(2T(r,h)+o(T(r,h)\right)
\end{align} 
for large $r$ possibly outside a set of finite logarithmic measure. Since $m\geq 2$,  from (\ref{EqT1.7}), we deduce that $m= 2$.

We now prove that $f$ has at most one pole. For the sake of simplicity, we suppose that $f$ has two poles at $a_1$ and $a_2$ of multiplicities $m_1$ and $m_2$ respectively. Assume that 
\begin{align}\label{EqT1.8}
f(z)=\frac{A(z)}{(z-a_1)^{m_1}(z-a_2)^{m_2}},
\end{align}
where $A(z)$ is a meromorphic function, which is analytic at $a_1$, $a_2$ and $A(a_1)\neq 0$ and $A(a_2)\neq 0$. Now from (\ref{EqT1.8}), we have 
\begin{align}\label{EqT1.9}
h(z)=\frac{A(z_1+z_2+\ldots+z_n)}{(z_1+z_2+\ldots+z_n-a_1)^{m_1}(z_1+z_2+\ldots+z_n-a_2)^{m_2}},
\end{align}
where $z=(z_1,z_2,\ldots,z_n)$. Now using (\ref{EqT1.9}) to (\ref{EqT1.0}), we get
\begin{align}\label{EqT1.10}
&\frac{B_i(z_1+\ldots+z_n)}{(z_1+\ldots+z_n-a_1)^{m_1+1}(z_1+\ldots+z_n-a_2)^{m_2+1}}\\&-\frac{A(z_1+\ldots+z_n)}{(z_1+\ldots+z_n-a_1)^{m_1}(z_1+\ldots+z_n-a_2)^{m_2}}\nonumber\\&=
a\frac{A(ng(z))}{(ng(z)-a_1)^{m_1}(ng(z)-a_2)^{m_2}}+c,\nonumber
\end{align}
where $i\in\{1,2,\ldots,n\}$ and
\begin{align*}
B_i(z_1+\ldots+z_n)&=(z_1+\ldots+z_n-a_1)(z_1+\ldots+z_n-a_2)\frac{\partial A(z_1+\ldots+z_n)}{\partial z_i}\\&-m_1(z_1+\ldots+z_n-a_2)A(z_1+\ldots+z_n)\\&-m_2(z_1+\ldots+z_n-a_1)A(z_1+\ldots+z_n)
\end{align*}
for $i\in\{1,2,\ldots,n\}$. Clearly $B_i(a_1)\neq 0$ and $B_i(a_2)\neq 0$ for $i\in\{1,2,\ldots,n\}$.

Let $\hat z\in\mathbb{C}^n$ be a zero of $z_1+\ldots+z_n-a_1$. Then from (\ref{EqT1.10}), we can say that $\hat z$ is a zero of $ng(z)-a_1$ or $ng(z)-a_2$. For the sake of simplicity, we may assume that $\hat z$ is a zero of $ng(z)-a_1$. If $\hat z$ is a simple zero of $ng(z)-a_1$, then from (\ref{EqT1.10}), we get a contradiction. Hence $\hat z$ is a $ng(z)-a_1$ of multiplicity atleast $2$. Since $\deg(g)=2$. it follows that $\hat z$ is a zero $ng(z)-a_1$ of multiplicity exactly $2$. Then we can write 
\begin{align}\label{EqT1.11}
ng(z)-a_1=\tilde a\left(z_1+\ldots+z_n-a_1\right)^2
\end{align}
where $\tilde a$ is a non-zero constant in $\mathbb{C}$.

Let $\tilde z\in\mathbb{C}^n$ be a zero of $z_1+\ldots+z_n-a_2$. Certainly from (\ref{EqT1.10}), we can say that $\tilde z$ is a zero of $ng(z)-a_2$. If $\hat z$ is a simple zero of $ng(z)-a_1$, then from (\ref{EqT1.10}), we get a contradiction. Hence $\hat z$ is a $ng(z)-a_2$ of multiplicity atleast $2$. Since $\deg(g)=2$. it follows that $\hat z$ is a zero $ng(z)-a_2$ of multiplicity exactly $2$. Then we can write 
\begin{align}\label{EqT1.12}
ng(z)-a_2=\tilde b\left(z_1+\ldots+z_n-a_2\right)^2,
\end{align}
where $\tilde b$ is a non-zero constant in $\mathbb{C}$.

Now from (\ref{EqT1.11}) and (\ref{EqT1.12}), we deduce that
\begin{align*}
\tilde a \left(z_1+\ldots+z_n-a_1\right)^2+a_1\equiv \tilde b\left(z_1+\ldots+z_n-a_2\right)^2+a_2,
\end{align*}
which is impossible here. We have thus proved that $f$ has at most one pole. Therefore, we can write
\begin{align}\label{EqT1.12a}
f(z)=h_1(z)h_2(z),
\end{align}
where $h_1(z)=\frac{1}{(z-w_0)^l}$, $z\in\mathbb{C}$, $l\in\mathbb{N}\cup\{0\}$ and $h_2(z)$ is an entire function in $\mathbb{C}$ with
$h_2(w_0)\neq 0$. Consequently
\begin{align}\label{EqT1.13}
h(z)=h_1(z_1+z_2+\ldots+z_n)h_2(z_1+z_2+\ldots+z_n),
\end{align}
where $z=(z_1,z_2,\ldots,z_n)$.

Now we consider the following two sub-cases.


\medskip
{\bf Sub-case 2.1.} Let $h_2$ be a transcendental entire function in $\mathbb{C}$. Using (\ref{EqT1.13}) to (\ref{EqT1.0}), we obtain
\begin{align}\label{EqT1.13a}
&\frac{\partial h_1(z_1+\ldots+z_n)}{\partial z_i}h_2(z_1+\ldots+z_n)+h_1(z_1+\ldots+z_n) \frac{\partial h_2(z_1+\ldots+z_n)}{\partial z_i}\\&=aG^g_{h_1}(z)G^g_{h_2}(z)+bh_1(z_1+\ldots+z_n)h_2(z_1+\ldots+z_n)+c,\nonumber
\end{align}
i.e.,
\begin{align}\label{EqT1.14}
&G^g_{h_2}(z)\\&=\frac{\frac{\partial h_1(z_1+\ldots+z_n)}{\partial z_i}h_2(z_1+\ldots+z_n)+h_1(z_1+\ldots+z_n) \frac{\partial h_2(z_1+\ldots+z_n)}{\partial z_i}}{G^g_{h_1}(z)}\nonumber\\&-\frac{bh_1(z_1+\ldots+z_n)h_2(z_1+\ldots+z_n)+c}{G^g_{h_1}(z)}\nonumber\\&=
h_2(z_1+\ldots+z_n)\left(\frac{\frac{\partial h_1(z_1+\ldots+z_n)}{\partial z_i}+h_1(z_1+\ldots+z_n) \frac{\frac{\partial h_2(z_1+\ldots+z_n)}{\partial z_i}}{h_2(z_1+\ldots+z_n)}-bh_1(z_1+\ldots+z_n)}{G^g_{h_1}(z)}\right)\nonumber\\&-\frac{c}{G^g_{h_1}(z)},\nonumber
\end{align}
where $z=(z_1,z_2,\ldots,z_n)$ and $i\in\{1,2,\ldots,n\}$. Note that $g(z)$ is a polynomial in $\mathbb{C}^n$ of degree $2$ and $h_1(z_1+\ldots+z_n)=\frac{1}{(z_1+\ldots+z_n-w_0)^l}$. Now applying Lemma \ref{L2} to (\ref{EqT1.14}), we deduce that
\begin{align}\label{EqT1.15}
T\left(r,G^g_{h_2}\right)&=m\left(r,G^g_{h_2}\right)\\&\leq m\left(r,h_2(z_1+\ldots+z_n)\right)+m\left(r,\frac{\frac{\partial h_2(z_1+\ldots+z_n)}{\partial z_i}}{h_2(z_1+\ldots+z_n)}\right)+O(\log r)\nonumber\\&
\leq T(r,h_2(z_1+\ldots+z_n))+o(T(r,h_2(z_1+\ldots+z_n)))
\end{align}
possibly outside a set of finite logarithmic measure. On the other hand by Lemma \ref{L1.1}, we have
\begin{align}\label{EqT1.16}
T\left(r,G^g_{h_2}\right)\geq (1-\varepsilon_3)T\left(r^2,h_2(z_1+\ldots+z_n)\right)
\end{align}
for any $\varepsilon_3>0$ and large $r$. We know that $T(r,h_2(z_1+\ldots+z_n))$ is an increasing convex function of $\log r$ so that$T(r,h_2(z_1+\ldots+z_n))/\log r$ is finally increasing and hence
\begin{align}\label{EqT1.17}
T(r,h_2(z_1+\ldots+z_n))\leq \frac{1}{2} T\left(r^2,h_2(z_1+\ldots+z_n)\right)
\end{align}
for sufficiently large $r$. Now from (\ref{EqT1.15})-(\ref{EqT1.17}), we obtain
\begin{align*}
T(r,h_2(z_1+\ldots+z_n))\leq \frac{1}{2}\frac{1}{1-\varepsilon_3}\left(T(r,h_2(z_1+\ldots+z_n))+o(T(r,h_2(z_1+\ldots+z_n))\right),
\end{align*}
which shows that $T(r,h_2(z_1+\ldots+z_n))=o(T(r,h_2(z_1+\ldots+z_n)))$. Therefore we get a contradiction.


\medskip
{\bf Sub-case 2.2.} Let $h_2$ be a polynomial in $\mathbb{C}$. Note that $g(z)$ is a polynomial in $\mathbb{C}^n$ of degree $2$ and $h_1(z_1+\ldots+z_n)=\frac{1}{(z_1+\ldots+z_n-w_0)^l}$. Now from (\ref{EqT1.13a}), we get
\begin{align}\label{EqT1.18}
&\frac{-lh_2(z_1+\ldots+z_n)}{(z_1+\ldots+z_n-w_0)^{l+1}}+\frac{1}{(z_1+\ldots+z_n-w_0)^l}\frac{\partial h_2(z_1+\ldots+z_n)}{\partial z_i}\\&=a\frac{h_2(ng(z))}{(ng(z)-w_0)^l}+b\frac{1}{(z_1+\ldots+z_n-w_0)^l}h_2(z_1+\ldots+z_n)+c,\nonumber
\end{align}
where $z=(z_1,z_2,\ldots,z_n)$ and $i\in\{1,2,\ldots,n\}$. Let $z^0=(z^0_1,z^0_2,\ldots,z^0_n)\in\mathbb{C}^n$ be a zero of $z_1+\ldots+z_n-w_0$. Since $h_2(w_0)\neq 0$, it follows that $h_2(z^0_1+\ldots+z^0_n)=h_2(w_0)\neq 0$ and $h_2(ng(z^0))=h_2(w_0)\neq 0$.
Now from (\ref{EqT1.18}), we can say that $\hat z$ is a zero of $ng(z)-w_0$. If $z^0$ is a simple zero of $ng(z)-w_0$, then from (\ref{EqT1.18}), we get a contradiction. Hence $z^0$ is a $ng(z)-a_1$ of multiplicity atleast $2$. Since $\deg(g)=2$. it follows that $z^0$ is a zero $ng(z)-w_0$ of multiplicity exactly $2$. Then we can write 
\begin{align}\label{EqT1.19}
ng(z)-w_0=\tilde a\left(z_1+\ldots+z_n-w_0\right)^2,
\end{align}
where $\tilde a$ is a non-zero constant in $\mathbb{C}$.

Now using (\ref{EqT1.19}) to (\ref{EqT1.18}) and then comparing the order of the pole at $z^0$ in the two sides of (\ref{EqT1.18}), we deduce that $l+1=2l$, i.e., $l=1$. Therefore from (\ref{EqT1.12a}), we get
\begin{align}\label{EqT1.19a}
f(z)=\frac{h_2(z)}{z-w_0},
\end{align}
where $z\in\mathbb{C}$. Also from (\ref{EqT1.18}), we have
\begin{align}\label{EqT1.20}
&\left(-lh_2(z_1+\ldots+z_n)+(z_1+\ldots+z_n-w_0)\frac{\partial h_2(z_1+\ldots+z_n)}{\partial z_i}\right)(ng(z)-w_0)\\&=ah_2(ng(z))(z_1+\ldots+z_n-w_0)^2+b(z_1+\ldots+z_n-w_0)h_2(z_1+\ldots+z_n)(ng(z)-w_0)\nonumber\\&+c(z_1+\ldots+z_n-w_0)(ng(z)-w_0),\nonumber
\end{align}
where $z=(z_1,z_2,\ldots,z_n)$ and $i\in\{1,2,\ldots,n\}$. 

Let $\deg(h_2)=d$. If $d\geq 2$, then it is easy to check that
the left hand side of (\ref{EqT1.20}) has degree $d+2$, while the right hand side has degree $2d + 2$, which is impossible. Thus, we have that $d = 0$ or $d = 1$. Therefore from (\ref{EqT1.19a}), we can write 
\begin{align}\label{EqT1.21}
f(z)=\beta+\frac{\alpha}{z-w_0},
\end{align}
for some constants $\alpha\neq 0, \beta$, where $z\in\mathbb{C}$. Now using (\ref{EqT1.21}) to (\ref{EqT1.0}), we get 
\begin{align*}
-\frac{\alpha}{(z_1+\ldots+z_n)^2}=\frac{a\alpha}{g(z)-w_0}+\frac{b\alpha}{z_1+\ldots+z_n-w_0}+c+(a+b)\beta,
\end{align*}
which shows that
\begin{align}\label{EqT1.22}
g(z)-w_0=\frac{a\alpha(z_1+\ldots+z_n-w_0)^2}{-\alpha-b\alpha(z_1+\ldots+z_n-w_0)-((a+b)\beta+c)(z_1+\ldots+z_n-w_0)^2},
\end{align}
where $z\in\mathbb{C}^n$.

Since $\deg(g)=2$, from (\ref{EqT1.22}), we conclude that $b\alpha=(a+b)\beta+c=0$, which shows that $b=0$ and $a\beta+c=0$. 
Consequently from (\ref{EqT1.22}), we have
\begin{align*}
ng(z)=w_0-a\left(z_1+\ldots+z_n-w_0\right)^2
\end{align*}

Therefore
\begin{align*}
f(z)=\beta+\frac{\alpha}{z-w_0}\quad \text{and}\quad ng(z)=w_0-a\left(z_1+\ldots+z_n-w_0\right)^2,
\end{align*}
$b=0$ and $\alpha\neq 0, \beta$ are constants such that $a\beta+c=0$.

Conversely, if $f$ and $g$ have the forms in Theorem \ref{t1.1} $(ii)$, it is easy to check that $f$ and $g$ satisfy the given equation. This proves Theorem \ref{t1.1} $(ii)$. The proof is thus complete.
\end{proof}

\begin{thebibliography}{99}
\bibitem[Hu-Li-]{Hu-Li-Yang-2003} {\sc P. C. Hu}, {\sc P. Li} and {\sc C. C. Yang}, Unicity of Meromorphic Mappings. Springer, New York (2003).
\end{thebibliography}
\end{document}
